% Pursuing Stacks (cotes 134-2 to 134-8): archive cote 134-5, site page 103 (PDF page 104), typescript page 290.
% typescript. Draft compared: ps4.tex:5297-5325.
% AI-transcribed by Claude Opus 5.5 (claude-opus-5-5, adaptive thinking, effort high), Message Batches API, batch ps2 (msgbatch_01NVwr1Yuo8qJf5rykeGBXPK), 2026-11-12; not yet checked by eye.
\scanpage{134-5/103}
\mspage{290}
as \xout{\illegible{1}} being generated by a contractibility structure.  Even when making
such an assumption,  it is not wholly clear yet that there is\hand{'}nt a good
notion of a morphism $f^{*}$ of contractibility structures,  without having
to assume there exists a left adjoint $f_{!}$ .  Maybe a little more ponde-%
ring on the situation would be useful.  If it should become clear that
(except for the obvious notion of \uline{equivalence} of asphericity structures)
there was \uline{not} any reasonable notion of a morphism of asphericity struc-%
tures,  this would probably mean that the notion of an asphericity
structure should not be viewed as a main structure type in homotopy
theory in it's own right,  but rather,  mainly,  as \deleted{an} \inserted{important} by\deleted{\unsure{e}}-product \deleted{\unsure{or|on}}
\deleted{an interesting aspect} of a cont\unsure{r}actibility s\unsure{t}ructure.  At any rate,
since section 81,  I feel that the main emphasis has \inserted{definitely} been shifting
towards contractibility structures,  which seems now \uline{the} main type of
structure dominating \xout{\unsure{the}} in the modelizing story,  as this \inserted{``story''} is gradually
emerging into light.

\margin{103}
